\section{Reproducibility}
\label{sec:reproducibility}

Every number in this paper is traceable to a released artifact through a recorded
chain of identifiers. This section documents that chain so a reader can verify a
value, re-derive a table, or run an equivalent experiment.

\subsection{Repository and Environment}

The implementation, the track configurations, the controllers, the evaluation
harness and the aggregated result artifacts are available in a public
repository~\cite{holodronecompetition2026}. The benchmark environment is
HoloOcean 2.3.0 with its Ocean world, Python 3.9.25 and the pinned dependency set
in \cmd{requirements.txt}; the learning extension uses a separate pinned set in
\cmd{requirements-rl.txt} installed in a distinct environment. Runs reported
here were executed on Windows 10 build 26200.

\subsection{Configuration-Driven Experiments}

A race is fully specified by one JSON file (Section~\ref{sec:configuration}):
world bounds and environment, the ordered gate sequence with centres, passage
directions and aperture sizes, the beacon network and its noise parameters, the
current profiles, the obstacle generation rules, the participants with their
vehicles, sensors, controllers and release delays, and the referee's validation
margins, penalties and scoring rules. A single entry point runs any of them, and
a controller is loaded by alias, by dotted module path or by file path plus class
name without modifying the package (Section~\ref{sec:controller}). Reproducing a
reported condition therefore requires the track file, the controller identifier
and the seed, all of which are recorded per run.

\subsection{Seeds and Determinism}

Beacon packet randomness is seeded by the run seed, the transmitter, the receiver
and the transmission index, so reception is independent of participant count and
of the order in which controllers are polled; a fleet run is reproducible for
this reason. Inter-vehicle packet loss and procedural course generation draw from
seeded generators. For the learning extension, seed bands are registered by role
--- training, validation, test and final holdout --- in a seed registry committed
to the repository, and each evaluation records both requested and completed
seeds.

\subsection{Recorded Artifacts}

Each run emits a line-delimited event log and a machine-readable summary, together
with the track, controller, seed and interface contract identifiers needed to
locate the corresponding result. Detailed provenance and artifact metadata are
retained in the repository release rather than repeated in the manuscript.

\subsection{Regenerating Tables and Figures}

All result tables and figures in this paper are \emph{post-processing} of
existing artifacts: no command listed below launches the simulator or modifies a
raw artifact.

\begin{lstlisting}[language=bash]
# Result tables and the penalty-identity check
python article/regenerate_tables.py --check

# Track layouts and the controller-comparison plot, from track JSON and artifacts
python article/figures/generate_figures.py

# Journal figures: perception panels from committed artifacts
python article_journal/scripts/make_perception_figure.py

# Manuscript
cd article_journal && latexmk -pdf -interaction=nonstopmode -halt-on-error main.tex
\end{lstlisting}

Running \cmd{regenerate\_tables.py} against the recorded result artifacts reproduces the
clean-track, current, fleet and coordination values of
Tables~\ref{tab:clean_tracks}, \ref{tab:currents} and \ref{tab:fleet} exactly, and
also verifies the penalty identity $T^{\mathrm{pen}}=T^{\mathrm{official}}+P$ on
every finished run. The two journal figure scripts write a provenance JSON
alongside each output recording their input files and asserting that no simulator
was launched.

\subsection{Re-Running Experiments}

Reproducing a run rather than a table requires HoloOcean and the recorded seed.
Every experimental command is stored verbatim in the per-run metadata; a clean
single-vehicle run has the form

\begin{lstlisting}[language=bash]
python -m marine_race_arena.scripts.run_marine_race \
  --track marine_race_arena/tracks/marine_race_horseshoe_bay.json \
  --benchmark-task clean_gate \
  --controller rule_gate_center_then_commit \
  --adapter holoocean --seed 0 \
  --duration 560.0 --current-profile none --official
\end{lstlisting}

and the current-free demonstration of Section~\ref{sec:eval_rq1} is driven by
committed scripts that pass \cmd{--current-profile none}. The control rate must
match the reported campaign: $30$\,Hz for the systematic benchmark evaluation
and $10$\,Hz for the current-free demonstration and learned-controller
validation.

\subsection{Updating the Learning Results}

The learning section is deliberately structured so that its validation can be
refreshed without changing the benchmark protocol. A new frozen checkpoint is
added by recording its selection decision, running the same three
seeds per official circuit, and regenerating Table~\ref{tab:learning} from the
per-episode result files. The methodology of Section~\ref{sec:learning}---the
27-D observation encoding, recurrent architecture, reward and evaluation
separation---is independent of any one checkpoint.
